\section{Detailed Trait-Conditioned Rule Induction and V30 Selection}
\label{app:trait_conditioned_details}

This appendix gives the implementation-level methodology underlying Section~\ref{sec:trait_conditioned_cktm}. It describes how trait-conditioned rules are induced from Q1--Q6 and how V30 information is selected for a particular reasoning step. The six C--K modes and the Q1--Q7 protocol are defined in Section~\ref{sec:ck_reasoning_modes} and Section~\ref{sec:participants_procedure}; Q7 is excluded from every procedure described below.

\subsection{Shared Data Protocol}

For participant $i$, let $\mathbf v_i=(v_{i1},\ldots,v_{i30})$ denote the V30 profile, where each item is scored on a six-point scale. All regression models use the standardized profile $\widetilde{\mathbf v}_i$. Resampling, permutation tests, and cross-validation are performed at the participant level, so that all transitions from the same designer remain in the same split. We use 1,000 bootstrap samples, 1,000 participant permutations, and five participant-grouped folds. This protocol prevents repeated transitions from one designer from inflating statistical support.

The procedures produce exploratory routing candidates rather than confirmed psychological mechanisms. A candidate is retained only when its effect is stable under resampling and improves held-out prediction; the associated permutation and multiple-testing diagnostics determine its evidence level. At generation time, candidates are used as soft preferences, never as deterministic state transitions or causal claims.

\subsection{Single-Trait Transition Rules}

For Q2, a trajectory $(z_{i1},\ldots,z_{iT_i})$ is expanded with \textsc{Start} and \textsc{Stop} boundaries. Repeated modes are retained. For each source mode $s$ and target mode $y$, we fit a one-versus-rest L1-regularized logistic regression:
\begin{equation}
\operatorname{logit} P(z_{i,t+1}=y\mid z_{it}=s)
= \alpha_{sy}+\boldsymbol\beta_{sy}^{\top}\widetilde{\mathbf v}_i
+\gamma_{1,sy}\phi_{it}+\gamma_{2,sy}\phi_{it}^{2},
\label{eq:app_univariate_transition}
\end{equation}
where $\phi_{it}=2t/T_i-1$ is the centered position of a transition. The L1 penalty is selected separately for each source mode by participant-grouped cross-validation over $C\in\{0.03,0.1,0.3,1.0,3.0\}$.

For every $(s,y,j)$ candidate, we record its standardized coefficient, participant support, bootstrap selection frequency, sign consistency, bootstrap interval, and held-out incremental log loss. A Q2 single-trait candidate enters the exploratory library when it has support from at least 20 participants, selection frequency of at least 0.90, sign consistency of at least 0.95, and $|\hat\beta_{sy,j}|\ge0.15$. Held-out comparisons add one trait dimension to a position-only baseline; participant-level permutation tests and Benjamini--Hochberg correction quantify the remaining uncertainty.

Coefficient polarity is compiled into an interpretable profile condition. A positive coefficient yields a high-pole trigger ($v_{ij}\ge4$); a negative coefficient yields a low-pole trigger ($v_{ij}\le3$). Both are rendered as a positive preference for the target operation under the corresponding profile condition. A negative coefficient therefore does not mean that the target operation is prohibited.

\subsection{Trait-Relation Transition Rules}

Single-trait effects cannot represent the possibility that two tendencies matter jointly. We first compute Spearman correlations for all unordered V30 pairs and retain stable pairs with $|\rho|\ge0.20$, bootstrap sign consistency of at least 0.90, and a bootstrap interval that excludes zero. The 12 strongest retained pairs define the relation-candidate set.

For each pair $(j,k)$, we construct the predicates in Table~\ref{tab:app_trait_predicates}. A predicate is considered only when both the predicate and its complement contain at least 20 participants.

\begin{table}[t]
\centering
\small
\caption{Trait-relation predicates used for pairwise rule induction.}
\label{tab:app_trait_predicates}
\begin{tabular}{ll}
\toprule
\textbf{Predicate} & \textbf{Activation condition} \\
\midrule
High co-activation & $v_{ij}\ge4\ \land\ v_{ik}\ge4$ \\
Low co-activation & $v_{ij}\le3\ \land\ v_{ik}\le3$ \\
Left tension & $v_{ij}\ge4\ \land\ v_{ik}\le3$ \\
Right tension & $v_{ij}\le3\ \land\ v_{ik}\ge4$ \\
Left dominance & $v_{ij}-v_{ik}\ge2$ \\
Right dominance & $v_{ik}-v_{ij}\ge2$ \\
Balance & $|v_{ij}-v_{ik}|\le1$ \\
\bottomrule
\end{tabular}
\end{table}

For a source--target pair and relation predicate $R_{jk}$, the baseline includes trajectory position and both trait main effects. The incremental model adds only the relation predicate:
\begin{equation}
\begin{aligned}
M_0:\quad \operatorname{logit}P(y=1) &= \alpha+\gamma_1\phi+\gamma_2\phi^2
+\beta_j\widetilde v_j+\beta_k\widetilde v_k,\\
M_1:\quad M_1 &= M_0+\theta R_{jk}(\mathbf v_i).
\end{aligned}
\label{eq:app_bivariate_transition}
\end{equation}
Thus, a retained relation rule reflects information beyond the two single-trait effects. We require positive residual support, bootstrap sign consistency of at least 0.90, improved held-out loss in at least three of five folds, and a favorable aggregate held-out loss. Negative relation effects are retained for analysis but are not compiled into generation preferences, because the routing language has no symmetric negative-action operator.

\subsection{Initial and Contextual Preference Rules}

Q1 records an initial mode rather than a transition. We therefore model $P(z_{i,\mathrm{initial}}=y\mid\widetilde{\mathbf v}_i)$ and compile an initial-orientation preference, without creating a fictitious source mode. Q3--Q6 likewise provide one mode choice in a stated dilemma. For each question $q$ and mode $y$, we fit
\begin{equation}
\operatorname{logit}P(Y_{iq}=y)
=\alpha_{qy}+\boldsymbol\beta_{qy}^{\top}\widetilde{\mathbf v}_i.
\label{eq:app_contextual_preference}
\end{equation}
The Q1 and Q3--Q6 models use a pre-specified L1 penalty of $C=0.3$. Their exploratory screen requires at least 15 supporting participants, selection frequency of at least 0.80, sign consistency of at least 0.90, and $|\hat\beta|\ge0.10$, together with favorable held-out diagnostics. Q1 produces initial-orientation candidates; Q3--Q6 candidates remain tied to their original conflict context and are activated only for semantically relevant situations.

\subsection{Rule Compilation and Runtime Routing}

Each retained candidate is compiled into a structured rule containing: (1) a V30 trigger, (2) a source operation or conflict context, (3) a preferred C--K operation, and (4) a factuality constraint. Statistical evidence is retained for audit but is not exposed to the language model. At generation time, routing is defined as
\begin{equation}
A_i=\operatorname{Route}(V_i,\mathcal R,X^\ast),
\label{eq:app_route}
\end{equation}
where $\mathcal R$ is the fixed exploratory candidate set. The router first removes rules whose profile trigger or semantic context does not apply, then ranks the remaining rules by trigger strength and held-out improvement. It retains at most one rule from each applicable rule type. If a single-trait and a trait-relation rule conflict for the same source operation, the relation rule is preferred because it controls for the two single-trait effects.

Rules do not edit the generated sequence. They only state preferred operations, and public task facts override them. Two constraints are attached to every applicable rule: a $C\rightarrow K$ operation must be expressed as future or conditional verification, and a $K\rightarrow K$ operation may infer consequences from known information but may not introduce a new proposal.

\subsection{Situation-Conditioned V30 Cue Selection}

The complete V30 profile is used internally, but serializing all 30
dimensions introduced redundant context in the development ablation. Before
generation, the selector therefore chooses four dimensions using their
relevance to the public Q7 situation and their salience for the target
designer:
\begin{equation}
\widetilde V_i(X^\ast)
=\operatorname{Top4}_{d\in\{1,\ldots,30\}}g(v_{id},X^\ast).
\label{eq:app_active4_v30}
\end{equation}
The selected dimensions are rendered as directional value cues without item
identifiers or numerical scores. Selection occurs once before prompting; it
does not depend on a predicted trajectory or an individual step.

The model then jointly produces aligned state--text units in one call:
\begin{equation}
\mathcal U_i\sim p_\psi(\mathcal U\mid
X^\ast,P_i,\widetilde V_i(X^\ast),H_i,R_i,E_i),
\qquad
\mathcal U_i=[(m_{it},y_{it})]_{t=1}^{T_i}.
\label{eq:app_one_call_generation}
\end{equation}
The sequence and response are deterministic projections of the same output.
No trajectory is fixed in a first continuation, and no second continuation
or step-conditioned V30 selector is used.
